\documentclass{article}
\usepackage{amsmath}
\usepackage{amsthm}
\usepackage{amssymb}
\usepackage{cleveref}
\usepackage{tikz}
\usepackage{float}

\setlength{\parindent}{0em}

\title{Problem Definition: Week 2}
\author{
    Ryan Batubara \\ 
    {\tt \small{rbatubara@ucsd.edu}} \\\And
    Ivy Hawks \\
    {\tt \small{mahawks@ucsd.edu}} \\\And
    Darren Ho \\
    {\tt \small{dah103@ucsd.edu}} \\\And
    Ciro Zhang \\
    {\tt \small{ciz001@ucsd.edu}} \\\And
    Shachar Lovett \\
    {\tt \small{slovett@ucsd.edu}} \\}
\normalsize

\begin{document}
\section*{Problem}
Consider a graph $G$ with edges $E$ and vertices $V$, and subgraphs $G_1, \ldots, G_k \subset G$. Let $G_i = (E_i, V_i)$ for these subgraphs. For now let us suppose the graph is unweighted and undirected. The subgraphs need not cover $G$ and can overlap. Each player $i$ knows $G$ and $G_i$, but not the other player's graphs. (This is akin to a map of a city, where each player is a different regional taxi company.) The goal is, given some starting vertex $v \in V$, and some end vertex, $w \in V$ to:
\begin{enumerate}
	\item  Identify if it is possible to construct a path $v, v_1, \dots, v_l, w$ such that every edge on the path is in $\bigcup_{i=0}^{k} E_i$. \label{one}
	\item  If such a path exists, find the minimum number of edges needed to get from $v$ to $w$, while using only edges belonging to $\bigcup_{i=0}^{k} E_i$. \label{two}
\end{enumerate}

\tikzset{
  point/.style={
    circle,
    fill,
    inner sep=1pt
  }
}

\begin{figure}[H]
\centering
\begin{tikzpicture}
  \draw (0,0) ellipse [x radius=6, y radius=2.5];
  \node at (0,2.1) {$G$};

  \draw plot [smooth cycle, tension=0.6] coordinates {(-2.2,-0.5) (-1.2,0.4)
	  (1.2,0.4) (1.8,-0.8) (0.2,-1.6) (-1.6,-1.7) };
  \node at (-1.1,-1.3) {$G_1$};

  \draw plot [smooth cycle, tension=1.1] coordinates {(-0.6,1.2) (0.8,1.6)
	  (2.1,0.6) (1.1,0) (-0.4,0.3)};
  \node at (0.7,1.2) {$G_2$};

  \draw plot [smooth cycle, tension=0.8] coordinates {(2.3,0.9) (4.6,1.2)
	  (4.4,-0.2) (3.2,-0.9) (1.0,-0.2)};
  \node at (4.2,0.8) {$G_3$};

  \node[point] (a1) at (-1.5,0.1) {};
  \node[point] (a2) at (-0.5,-0.6) {};
  \node[point] (a3) at (0.8,-0.4) {};
  \draw[dotted] (a1) -- (a2) -- (a3);

  \node[point] (b1) at (-0.2,0.9) {};
  \node[point] (b2) at (0.8,0.7) {};
  \node[point] (b3) at (1.4,0.2) {};
  \draw[dotted] (b1) -- (b2) -- (b3);

  \node[point] (c1) at (2.6,0.5) {};
  \node[point] (c2) at (3.8,0.6) {};
  \node[point] (c3) at (3.4,-0.3) {};
  \draw[dotted] (c1) -- (c2) -- (c3);

  \node[point] (d1) at (-5,0) {};
  \node[point] (d2) at (-3,-0.7) {};
  \node[point] (d3) at (-2.4,1) {};
  \node[point] (d4) at (-4,0.4) {};
  \draw[dotted] (d1) -- (d2) -- (d4);
  \draw[dotted] (d2) -- (d3) -- (b1);

  \node[below left=0.1pt , outer sep=2pt] at (0.8,-0.4) {$v$};
  \node[above right=1pt , outer sep=2pt] at (2.6,0.5) {$w$};

  \draw[thick] (c1) -- (b3) -- (a3);
\end{tikzpicture}
\caption{A valid path from $v$ to $w$ in bold. Note that $G$ may contain nodes
not belonging to any player.}
\end{figure}

\newpage
\section*{Lower bound sketches and ideas}
For \cref{one}, we think this would require computing $\binom{k}{2}$ pairwise set disjointnes in the worst case, to determine where the intersections of each $G_i$ are, at which point each player could on their own determine the possible
paths between each of the subgraphs that intersect with their own. 

\begin{figure}[H]
\centering
\begin{tikzpicture}
  \draw (0,0) ellipse [x radius=6, y radius=2.5];
  \node at (0,2.1) {$G$};

  \draw plot [smooth cycle, tension=0.9] coordinates {(-4.2,0.5) (-3.2,1.4) (-1.8,0.8) (-2.2,-0.6) (-3.6,-0.7)};
  \node at (-3.1,1) {$G_1$};

  \draw plot [smooth cycle, tension=1.1] coordinates {(-0.6,1.2) (0.8,1.6) (2.1,0.6) (1.1,-0.2) (-0.4,0.3)};
  \node at (0.7,1.2) {$G_2$};

  \draw plot [smooth cycle, tension=0.8] coordinates {(2.3,0.9) (4.6,1.2) (4.4,-0.2) (3.2,-0.9) (2.0,-0.2)};
  \node at (4.2,0.8) {$G_3$};

  \node[point] (a1) at (-3.6,0.2) {};
  \node[point] (a2) at (-2.6,0.6) {};
  \node[point] (a3) at (-2.8,-0.4) {};
  \draw[dotted] (a1) -- (a2) -- (a3);

  \node[point] (b1) at (-0.2,0.9) {};
  \node[point] (b2) at (0.8,0.7) {};
  \node[point] (b3) at (1.4,0.2) {};
  \draw[dotted] (b1) -- (b2) -- (b3);

  \node[point] (c1) at (2.6,0.5) {};
  \node[point] (c2) at (3.8,0.6) {};
  \node[point] (c3) at (3.4,-0.3) {};
  \draw[dotted] (c1) -- (c2) -- (c3);
\end{tikzpicture}
\caption{Worst case for \cref{one}: Disjoint Subsets}
\end{figure}


For \cref{two}, we would be required to compute $\binom{k}{2}$ pairwise
intersections, since in order to minimize the path we would need to know all
possible paths that could connect $v$ and $w$. To find the shortest path, each
player works backwards and sends the distances through their internal vertices
to the players they intersect with, and each player will attempt to identify if
a longer path can be dropped whenever they receive this list.
\end{document}
